\textbf{COLUNI 2021 -- Questão 04}\par\medskip No livro \textit{A menina que fez a América}, de Ilka B. Laurito, a preocupação com o trabalho e/ou o cuidado com a família são características importantes dos imigrantes italianos. Identifique os excertos abaixo que melhor justificam essa afirmativa.

\par

I. “Sua mãe deve ter juntado algum dinheiro, mas a vida no Brasil é difícil. Uma passagem de navio custa muito caro. Além disso, a terceira classe dos emigrantes.” (LAURITO, Ilka Brunhilde. \textbf{A Menina que fez a América}. São Paulo, FTD, 2019. p. 83) II. “Em fins do século XIX, a vida estava difícil para os trabalhadores dos campos da Itália e de outras partes da Europa. Faltava serviço e, quando havia, era mal pago.” (LAURITO, Ilka Brunhilde. \textbf{A Menina que fez a América}. São Paulo, FTD, 2019. p. 08) III. “Comecei a choramingar baixinho, depois alto, e se Bettina tivesse demorado um pouco mais eu teria feito transbordar a fonte da Acquanova com meu pranto e, no dia seguinte, as mulheres lavariam a roupa com as minhas lágrimas...” (LAURITO, Ilka Brunhilde. \textbf{A Menina que fez a América}. São Paulo, FTD, 2019. p. 70) IV. “Continuamos a subir, com os vagões gemendo atrás da locomotiva fumegante, e eu passei a imaginar que a cidade fosse do tamanho de uma noz, daquelas que balançavam nos mais altos galhos das nogueiras da Calábria.” (LAURITO, Ilka Brunhilde. \textbf{A Menina que fez a América}. São Paulo, FTD, 2019. p. 92)

\par

Marque a alternativa CORRETA.

\par

a) II e IV. b) I e III. c) I e II. d) III e IV.
